% Generated draft: requires mathematical and source review.
\documentclass[11pt]{article}
\usepackage[margin=25mm]{geometry}
\usepackage{fontspec}
\setmainfont{Latin Modern Roman}
\usepackage{amsmath,amssymb,mathtools,unicode-math}
\setmathfont{Latin Modern Math}
\usepackage{enumitem,needspace,xurl,hyperref}
\setlength{\parindent}{0pt}
\setlength{\parskip}{5pt}
\setlength{\emergencystretch}{3em}
\newcommand{\N}{\mathbb{N}}
\newcommand{\Z}{\mathbb{Z}}
\newcommand{\Q}{\mathbb{Q}}
\newcommand{\R}{\mathbb{R}}
\newcommand{\C}{\mathbb{C}}
\newcommand{\F}{\mathbb{F}}
\newcommand{\A}{\mathbb{A}}
\newcommand{\PP}{\mathbb{P}}
\newcommand{\E}{\mathbb{E}}
\newcommand{\eps}{\varepsilon}
\newcommand{\dd}{\,\mathrm{d}}
\begin{document}
\section*{UnsolvedMath ID: 800381}
\textbf{Source problem code:} catalogue entry
\subsection{Definitions and mathematical statement}
Let \(\Q^*=\Q\setminus\{0\}\). For each \(b,c\in\Q^*\), define the degree-two rational map
\[
f_{b,c}\colon\PP^1_{\Q}\longrightarrow\PP^1_{\Q},\qquad [X:Y]\longmapsto[bX^2+cY^2:XY].
\]
On the affine locus where \(x\ne0\), this is \(f_{b,c}(x)=bx+cx^{-1}\). The projective extension gives \(f_{b,c}(0)=\infty\) and \(f_{b,c}(\infty)=\infty\). Thus the affine formula's boundary points are included in the projective formulation.

Write \(f^{\circ0}\) for the identity and \(f^{\circ n}\) for the \(n\)-fold iterate. A point \(P\in\PP^1(\Q)\) is periodic if \(f^{\circ n}(P)=P\) for some integer \(n\ge1\); its exact period is the least such \(n\). Define
\[
\operatorname{PrePer}(f,\PP^1(\Q))
=\{P\in\PP^1(\Q): f^{\circ m}(P)=f^{\circ n}(P)\text{ for some integers }0\le m<n\}.
\]
The supplied conjecture, attributed to Manes (2008), asserts that for every \(b,c\in\Q^*\), every periodic point in \(\PP^1(\Q)\) has exact period in \(\{1,2,4\}\), and that
\[
\sup_{b,c\in\Q^*}\#\operatorname{PrePer}(f_{b,c},\PP^1(\Q))=12.
\]
The latter assertion means that every such map has at most twelve rational preperiodic points and that some such map has exactly twelve. These are conjectural assertions in the supplied record. The supplied Manes reference lacks full bibliographic details.
\subsection{Short English statement}
For every pair of nonzero rational parameters \(b,c\), the conjecture predicts that rational periodic points of \(bx+c/x\), viewed on the projective line, have exact period \(1\), \(2\), or \(4\). It also predicts a sharp uniform maximum of twelve rational preperiodic points, including any points at zero and infinity.
\subsection{Sources}
\begin{itemize}
\item Local UnsolvedMath registry, permanent ID \(800381\), source problem code \texttt{catalogue entry}, titled \emph{Open Problems in Arithmetic Dynamics --- Conjecture 4}. The registry supplies the status \emph{open}; this status has not been independently verified.
\item Benedetto, DeMarco, Ingram, Jones et al., \emph{Current Trends and Open Problems in Arithmetic Dynamics} (2018), Conjecture 4. Supplied source URL: \url{https://arxiv.org/abs/1806.04980}. The supplied provenance records access on 2026-10-16 and extraction by \texttt{arxiv\_latex}; these are supplied provenance claims, not a claim of independent access or verification.
\item The supplied statement attributes the conjecture to Manes (2008), with reference label \texttt{MR2407816}. No full bibliographic entry for this reference is supplied.
\end{itemize}
\end{document}
